% SOURCE_STATEMENT: "For edge-intensified tropical convection, determine (1) what sets the number and separation distance of convergence lines, including without a sea-surface-temperature gradient; (2) why convection does not occur over the highest sea-surface temperature and whether a predictive model of edge intensification can be derived; and (3) how rotation changes these effects."
% Preparation and final self-audit: CONTRIBUTING.md prompts C and D; source/audit artifacts retained in the local physics_statements_2026_10_11 batch.
% !TeX program = lualatex
\documentclass[11pt]{article}
\usepackage[margin=25mm]{geometry}
\usepackage{fontspec}
\setmainfont{Latin Modern Roman}
\usepackage{amsmath,amssymb,mathtools}
\usepackage{unicode-math}
\setmathfont{Latin Modern Math}
\usepackage{xurl,hyperref}
\hypersetup{colorlinks=true,urlcolor=blue}
\setcounter{secnumdepth}{0}
\setlength{\parindent}{0pt}
\setlength{\parskip}{5pt}
\setlength{\emergencystretch}{3em}
\newcommand{\sref}[2]{\hyperlink{source:#1:#2}{[S#2]}}
\newcommand{\cataloguescope}[1]{\paragraph{Recorded target.}#1}
\newcommand{\R}{\mathbb R}
\newcommand{\C}{\mathbb C}
\newcommand{\N}{\mathbb N}
\newcommand{\Z}{\mathbb Z}
\newcommand{\Q}{\mathbb Q}
\newcommand{\F}{\mathbb F}
\newcommand{\eps}{\varepsilon}
\begin{document}
\section{Mechanisms and Scales of Tropical-Cyclone Edge Intensification}\label{statement}

\subsection{Definitions and mathematical statement}

\paragraph{Shared notation.}$\mathbb N=\{1,2,\ldots\}$; the standard number fields have their usual meanings. All additional conventions are stated locally.


Edge intensification refers to enhanced convection near the edge of a warm region rather than its warmest interior. Convergence lines collect horizontal flow feeding rising air. Their number and spacing are observables of the convection pattern; rotation introduces Coriolis effects. Sea-surface-temperature gradients are possible forcing rather than a stipulated necessary cause.

\paragraph{Catalogue wording retained for traceability.}
For edge-intensified tropical convection, determine (1) what sets the number and separation distance of convergence lines, including without a sea-surface-temperature gradient; (2) why convection does not occur over the highest sea-surface temperature and whether a predictive model of edge intensification can be derived; and (3) how rotation changes these effects.
\subsection{Short English statement}
Explain convergence-line spacing, edge-focused convection and the influence of rotation, including settings without a sea-surface-temperature gradient.
\subsection{Sources}
\begin{itemize}
\item[\hypertarget{source:statement:1}{S1}] ULAM.AI and the UnsolvedMath Contributors, UnsolvedMath, catalogue entry. Catalogue curation under CC BY 4.0; definitions and English exposition added locally.
\url{https://huggingface.co/datasets/ulamai/UnsolvedMath}.
\item[\hypertarget{source:statement:2}{S2}] Model Hierarchies in Atmosphere, Ocean, and Climate Sciences (2024), Oberwolfach Reports, DOI 10.4171/owr/2024/31, pp. 1731--1786. Original report and mathematical context. Incorporated report material remains under the source CC BY-SA 4.0 terms unless otherwise noted.
\url{https://doi.org/10.4171/owr/2024/31}.
\end{itemize}
\label{end:statement}
\end{document}
